J'ai une activit� r�guli�re d'encadrement doctoral (dont je ne sais si elle doit plut�t figurer dans la section {\sl Recherche} ou dans la section {\sl Enseignement}). Depuis 2010, j'ai (co-)encadr� une douzaine d'�tudiants dont la liste est donn�e en section \ref{sec:theses} et dont j'ai mentionn� les travaux au chapitre pr�c�dent. La plupart de ces encadrements se sont fait en collaboration avec un ou plusieurs coll�gues, soit parce que le sujet le requiert (dans le cadre d'une th�se � l'interface entre statistique et biologie par exemple), soit pour permettre � des coll�gues plus jeunes d'avoir une premi�re exp�rience d'encadrement. 

L'encadrement de th�se constitue une part centrale dans mon activit� de recherche et les �tudiants que j'ai la chance d'encadrer font partie du premier cercle de mes collaborateurs. J'ai eu la chance depuis plusieurs ann�es de travailler avec des �tudiants tr�s bien form�s avec qui la collaboration scientifique a �t� riche et efficace. Ma propre production scientifique b�n�ficie largement de ces exp�riences. J'esp�re contribuer ainsi � former des jeunes chercheurs tr�s solides en math�matiques et int�ress�s aux sciences de la vie, qui pourront ensuite contribuer de fa�on significative � la recherche � l'interface entre statistique et biologie. La majorit� d'entre eux occupent aujourd'hui des postes acad�miques : ma�tre de conf�rences, charg� de recherche, ing�nieur de recherche.
